\chapter{过渡金属：颜色、磁性与催化}

\begin{kaichang}
请在身边找齐这四样东西：一枚生了锈的铁钉，一截从废旧电线里剥出来的红铜丝，一把亮闪闪的不锈钢勺子，一块冰箱贴上的小黑磁铁。

现在考考你的眼力。铁钉和不锈钢勺——按周期表（第 10 章那张地图）来说，它们的主要成分住在同一个街区，凭什么一个锈迹斑斑，一个用了十年还光亮如新？铜丝是红棕色的，而大多数金属都是银白色，铜凭什么搞特殊？还有那块磁铁：世界上一百多种元素，能被磁铁狠狠吸住的常见金属，翻来覆去就是铁、钴、镍这三兄弟——它们仨在周期表上恰好是邻居。这是巧合吗？

再补一个悬念：化学实验室里有一种蓝色晶体，叫胆矾，丢进水里就化开成一杯漂亮的蓝宝石色溶液。它的成分不过是铜、硫、氧和水——硫酸盐明明大多是无色的，这杯蓝色到底是谁染上去的？

这四样东西的答案，都藏在周期表中间那片我们还没细逛过的街区里。这一章，我们挨家挨户去敲门。
\end{kaichang}

\section{中间街区的三样特产}

第 10 章把过渡金属叫作周期表中央“一大片矮房子”：铁、铜、锌、钛、铬、锰、钴、镍都住在这里，往下还有银、金、铂、汞这些熟面孔。这一街区有三个特产，是左右两边街区都拿不出来的。

第一样特产是\textbf{颜色}。主族元素的世界相当朴素：食盐、白糖、石英、纯碱，几乎都是白色或无色的。而过渡金属这边像个颜料铺：含铜离子的溶液是蓝色的，高锰酸钾（消毒用的紫药水原料）溶液是深紫色的，含铬离子的溶液是绿色的，铁锈是红棕色的。宝石商最懂这一点——红宝石和蓝宝石的骨架其实是同一种无色矿物（刚玉，成分是氧化铝），只不过红宝石里混进了万分之几的铬原子，蓝宝石里混进了铁和钛。一点点过渡金属杂质，就把无色的石头变成了国宝。

第二样特产是\textbf{磁性}。把一把元素样品挨个凑到磁铁跟前，绝大多数毫无反应，只有铁、钴、镍会被牢牢吸住，而且它们自己也能被做成磁铁。你的冰箱贴、耳机里的喇叭、电动车的电机，全靠这三兄弟。

第三样特产是\textbf{催化}。催化剂的“催化”两个字，意思是让反应走得更快，自己却不在最终产物里。化肥厂把空气中的氮气变成氨，用的是铁做的催化剂；汽车排气管里那个贵重的“三元催化器”，用的是铂、钯、铑；你实验室里往双氧水里丢的那勺黑色二氧化锰粉末，主角是锰。翻开工农业的催化剂名录，过渡金属和他们的化合物占了大半壁江山。

三样特产，三条线索。这一章的任务是让你看到：它们其实是从同一个根源上长出来的。

\section{多出来的十个格子：d 电子登场}

回到周期表的形状问题。第 10 章说过，第一行 2 格、第二三行各 8 格，是因为电子按“楼层”依次入住，第一层住 2 个、第二层住 8 个。可到了第四行，表格突然从 8 格撑宽到 18 格——多出来的 10 个格子，就是第一排过渡金属，从钪一直排到锌。

这 10 个格子对应电子入住的一类新“房间”：d 房间。同一层里，d 房间一共 5 间，每间住 2 个电子，正好 10 个名额。这里有个关键细节，第 10 章的挑战层提过一句，现在要正式用了：电子入住的顺序有点别扭——第四层的 s 房间比第三层的 d 房间“便宜”（能量更低），所以第四行的元素先把 4s 住上，再回过头来慢慢填 3d。从钪到锌，一个一个新增的电子，进的正是这个“回马枪”式的 3d 房间。

为什么这个别扭如此重要？因为它造成了一个独一无二的局面：\textbf{3d 电子和 4s 电子的能量几乎一样高}。主族元素的世界里，电子楼层的能量差得很开，内层电子被牢牢按住，永远轮不到出场，只有最外层那一两个电子参与化学。而在过渡金属这里，3d 和 4s 住得太近了，近到化学反应的能量足以把它们挨个请出来。

直接后果就是：\textbf{同一元素可以交出不同数目的电子，呈现好几种氧化态}。铁可以交出 2 个电子变成 \chem{Fe^{2+}}，也可以再交一个变成 \chem{Fe^{3+}}；铜有 \chem{Cu^{+}} 和 \chem{Cu^{2+}}；锰最夸张，从 $+2$ 到 $+7$ 一口气排出七八种价态——紫药水里的高锰酸根 \chem{MnO_4^{-}} 就是 $+7$ 价。主族的钠永远只交出 1 个电子，一辈子一个价；而锰一个人就能演出七八种角色。价态多变，正是颜色多变、催化神通的前奏。

\section{颜色从哪里来：被吃掉的那部分光}

先想清楚一个基本问题：我们说的“颜色”到底是什么？白光——阳光、灯光——是各种颜色光的混合。一样物体呈现什么颜色，取决于它把白光里的哪部分\textbf{吃掉}（吸收）了。一片叶子把红光和蓝紫光吃掉，剩下绿光弹回你的眼睛，叶子就是绿的。所以问“硫酸铜溶液为什么是蓝色”，等于问：“是谁吃掉了橙黄色的光？”

答案落在 d 电子身上。过渡金属离子的 3d 电子住进了 5 间 d 房间，这些房间的能量并不完全相同，电子可以从低的房间跃迁到高的房间——只要外来的一份光能量恰好等于两级之差，这份光就被吃掉。关键在于：\textbf{d 房间之间的能量差，恰好落在可见光的能量范围内}。不多不少，正好。于是过渡金属离子吃进一部分可见光，把剩下的颜色还给我们。

对比一下就明白了。钠离子、钙离子这些主族离子，电子也住在房间里，可它们的空房间太远，能量差落在紫外线范围——人眼看不见紫外线，吃掉多少都无所谓，溶液就是无色的。食盐水透明，不是因为它“什么都不吸收”，而是因为它吸收的光你看不见。

还有一个更精彩的细节：d 房间之间的能量差不是固定的，它会被周围的东西\textbf{拉扯改变}。这就引出了本章最重要的概念——配位。

\section{配位：金属离子身边围着一圈“房客”}

\chem{Cu^{2+}} 在水里从来不是光杆司令。水分子氧原子上有富余的电子对，会被带正电的铜离子吸引，四到六个水分子围拢过来，把孤对电子“借”给铜离子使用，形成一圈稳定的包围。这样的整体叫\textbf{配位化合物}（或配合物）：中心是金属离子，围在身边、出借电子对的分子或离子叫\textbf{配体}。硫酸铜溶液的蓝色，严格说不属于 \chem{Cu^{2+}} 本身，而属于“铜离子加六个水配体”这个组合。

证据随手可做：无水硫酸铜（把水配体全部赶跑）是灰白色的粉末，一旦遇水，水分子围拢上去，蓝色立刻现身。换一种配体，颜色就变——往蓝色铜溶液里加浓氨水，氨分子把水分子挤走、自己围上去，溶液会转成深得多的蓝紫色（这个实验产生刺激性气体，请看老师演示或视频，不要在家做）。同一个铜离子，只因身边的房客换了，d 房间被拉扯出的能量差变了，吃掉的光变了，颜色就变了。

配体不同、颜色不同，这条规律解释了一大票现象：红宝石的铬离子嵌在氧化铝晶格里，被氧原子包围，是红色；同样是铬，换到祖母绿的绿柱石晶格里，周围的包围方式变了，就成了绿色。宝石的颜色，本质上是配体替铬离子“调”出来的。

\section{磁性：排队站好的亿万根小磁针}

颜色讲完了，磁性呢？根源还在 d 电子，不过这次的关键词是\textbf{自旋}。

第 9 章讲过，电子有一种叫自旋的属性，直觉上可以想象成每颗电子都是一根微小到极点的磁针（这是帮助思考的模型，不是电子真的在转）。当电子成对住进同一房间时，两根小磁针方向相反，磁性恰好抵消。可过渡金属的 3d 房间常常是“半住满”状态——比如 \chem{Fe^{3+}} 的 5 个 d 电子，按洪特规则（第 9 章）一人独占一间房，5 根小磁针全部指向同一方向，谁也抵消不了谁。这样的离子就是个微型磁铁。

但光有这种小磁铁还不够。把一块普通铁块放在桌上，它并不吸铁钉。为什么？因为铁块内部划分成无数个微小区域，叫\textbf{磁畴}：每个磁畴内部，亿万根小磁针方向一致；但磁畴与磁畴之间指向乱七八糟，互相抵消。铁、钴、镍的真正特殊之处在于：它们原子间的相互作用能让相邻小磁针\textbf{对齐的成本特别低}——拿一块外磁铁靠近，磁畴像听到口令一样集体转向、整齐排列，整块铁立刻变成强磁体，这就是“磁化”。铜、铝、锌的 d 电子要么成对抵消，要么原子间没有这种齐心协力的对齐机制，于是与磁性无缘。三兄弟的邻居关系不是巧合：3d 房间住的人数相近，才让它们的原子恰好满足“小磁针多、又容易对齐”这两个苛刻条件。

\section{铁锈、不锈钢和一场慢速的电化学火灾}

现在解决开场的第一个悬念：为什么铁会生锈，不锈钢不会？

先看铁锈。把铁钉放在干燥空气里，几年都安然无恙；泡进煮过又密封的水里，也能撑很久；可一旦水和氧气同时到场，红棕色的锈就爬出来了。生锈的实质是铁把电子交给氧气：铁原子先交出电子变成 \chem{Fe^{2+}}，电子沿着金属和水膜传给氧，再经过一系列转变，最终形成含水的氧化铁——主要成分可以写作 \chem{Fe_2O_3}：
\[\chem{4Fe} + \chem{3O_2} \rightarrow \chem{2Fe_2O_3}\]
真实的铁锈还裹挟着数量不定的水分子，质地疏松，像一块吸满水的海绵。疏松是关键败笔：它挡不住里面的铁继续接触水和氧，于是锈一层一层往里啃，直到整根铁钉报废。

不锈钢的主角还是铁，但混进了约 $12\%$ 以上的铬。铬和氧的亲和力更强，抢先在表面生成一层氧化铬。和铁锈相反，这层氧化铬\textbf{致密、牢固、透明}，厚度只有几纳米，却把空气和水挡得严严实实。更妙的是它有自愈能力：勺子被划伤，露出的铬立刻再氧化，薄膜自动补上。你看不见这层膜，但它每时每刻都在站岗。所以“不锈钢不生锈”准确的说法是：它锈得太快、锈得太好，锈出一层完美的盔甲，然后就不再锈了。

人类防腐的全部招数，都是沿着这条思路展开的：刷漆、涂油，是物理上隔绝水和氧；镀锌（白铁皮）多了一层心机——锌比铁更容易交电子，即使镀层破损，锌也抢着先被腐蚀，替铁“牺牲”掉，这叫牺牲阳极保护法；轮船外壳和跨海大桥的钢桩上，就焊着成块的锌或镁，定期更换。

\begin{shiyan}
\textbf{铁钉生锈对比实验}（在家可做，全程安全）

材料：四枚干净铁钉、四个透明杯子、自来水、食盐、食用油、纸巾。

步骤：① 一号杯放铁钉，杯口敞开，只加干燥的空气（杯底垫干纸巾）；② 二号杯加自来水没过铁钉一半，让铁钉一半在水里、一半在空气里；③ 三号杯加溶了两勺食盐的水，同样没过铁钉一半；④ 四号杯加自来水完全没过铁钉，再倒一层食用油封住水面。

观察点：接下来三到七天，每天看一次并记录。哪枚锈得最快？锈斑先出现在什么位置（水里还是水面交界处）？你会发现：干湿交界的二号锈得明显，盐水里的三号锈得更快，被油封住的四号和最干燥的一号几乎不锈。

结论等你亲手得出：生锈需要水和氧\textbf{同时}到场，而盐会让这场“慢速火灾”烧得更旺（海水里的船为什么这么容易锈，你现在知道了）。
\end{shiyan}

\section{催化剂：给反应修一条穿山隧道}

第三样特产轮到催化了。先用第 27 章还会细讲的能量语言把话说准：一个反应能不能发生、能进行到什么程度，由能量规则（热力学）决定；而反应有多快，由另一套规则（动力学）决定——关键障碍是\textbf{活化能}，反应物必须先翻过一座能量山，才能变成产物。山越高，能翻过去的分子越少，反应越慢。

催化剂做的事只有一件：\textbf{给反应换一条山更低的路}。它不推反应物一把，不供给能量，也不能把翻不过山的反应变成能翻过去——一座根本翻不过去的山，修隧道也无济于事；它只是让本来能发生、但慢得没有实用价值的反应，快到够用。山两边的起点和终点没变，所以催化剂不改变反应放多少热，也不改变最终平衡停在什么位置。

过渡金属凭什么特别会修隧道？答案还是 d 电子和多变价态。以固体表面催化为例：化肥厂把氮气和氢气变成氨，
\[\chem{N_2} + \chem{3H_2} \rightleftharpoons \chem{2NH_3}\]
这条路最大的障碍是氮气分子里那根异常牢固的三键。在铁催化剂表面，\chem{N_2} 分子先被铁原子表面的 d 电子“抓住”（吸附），氮原子一个个趴在铁表面，被多只手同时拉住，三键被削弱、拆散；氢气同样被拆开成氢原子；然后氮、氢原子在表面这个“工作台”上逐个接头，拼成氨分子，最后脱手离开。表面把分子拆开、按住、排队、撮合——一条原本不存在的低能量路径就这样开张了。反应结束，铁表面恢复原样，等待下一批客人。铂、钯、铑在汽车三元催化器里干的是同类工作：把尾气里的一氧化碳、氮氧化物、未烧净的碳氢化合物抓上表面，撮合它们变成二氧化碳、氮气和水再放走。你的汽车排气管里那小小一罐，是世界上最昂贵的几克金属在上班。

\begin{qiaoliang}
“快一点”到底是多快？活化能的影响是指数级的：在其他条件相近时，反应速率大致正比于 $e^{-E_a/(RT)}$，其中 $E_a$ 是活化能，$T$ 是温度，$R$ 是常数。在室温附近，$RT$ 大约是 \ud{2.5}{kJ\,mol^{-1}}。假如催化剂把活化能降低 \ud{20}{kJ\,mol^{-1}}（这在催化里只是中规中矩的成绩），速率将增大约
\[e^{20/2.5} = e^{8} \approx 3000\ \text{倍}\]
降低 \ud{40}{kJ\,mol^{-1}}，则是 $e^{16}$，约一千万倍。指数的意思是：活化能每降一小截，速率不是加几倍，而是\textbf{乘}几倍。这就是为什么“换个催化剂”常常是化工厂里比“升温加压”更省钱的选择——而生命体里的酶（蛋白质做的催化剂，第二册的主角）把这场指数游戏玩到了登峰造极。
\end{qiaoliang}

\begin{wuqu}
“催化剂参与了反应，所以会被慢慢消耗掉；用得越久，效果越差，说明它被用光了。”——前半句对，后半句错。催化剂确实深度参与反应：抓分子、拆键、撮合，一样没少干；但它在每一轮循环结束时\textbf{恢复原状}，理论上不被消耗。工业催化剂用久了效果变差，真实原因多半是“中毒”和“磨损”：杂质（比如硫）牢牢粘在表面，霸占了工作台；或者高温让表面烧结、面积缩小。这不是催化剂被用掉，而是工作台被弄脏、被烧坏了。反例就在眼前：实验室里那勺二氧化锰催化双氧水分解，反应完过滤、干燥，称一称，质量分毫不少，下回接着用。
\end{wuqu}

\section{科学家怎样知道：一场关于颜色的百年争论}

\begin{zhengju}
19 世纪末，化学家手里有一批让人头疼的化合物：把氨通进钴盐溶液，能得到 \chem{CoCl_3 \cdot 6NH_3}、\chem{CoCl_3 \cdot 5NH_3}、\chem{CoCl_3 \cdot 4NH_3} 等一系列物质——钴一样多，氨一个个减少，颜色却从黄到紫到绿，各不相同。按当时的理论，原子之间靠一根根“键”连成链条，可没人画得出这些链：氨分子凑在氯化钴旁边，到底连在哪？

瑞士化学家维尔纳在 1893 年提出一个当时被视为异端的模型：钴离子在正中央，氨分子和氯离子\textbf{围成一圈}直接贴着钴——这一圈叫内界；多余的氯离子待在外圈，只起平衡电荷的作用。这个模型做出了可以被推翻的预言：如果氯在内圈，它就和钴绑在一起，掉进水里也不会以自由氯离子出现；在外圈的氯才会。于是实验变得干脆利落：溶解样品，加入硝酸银，看能立刻沉淀出多少氯化银。结果 \chem{CoCl_3 \cdot 6NH_3} 沉淀出 3 份氯，\chem{CoCl_3 \cdot 5NH_3} 只沉淀出 2 份，\chem{CoCl_3 \cdot 4NH_3} 只沉淀出 1 份——和预言严丝合缝。反对者也曾提出替代解释（比如氨先抱成团再挂在链上），但每种新解释都被新的电导率、沉淀计量实验逐一排除。更硬的证据还在后面：维尔纳预言某些配合物应该有像左右手一样不能重合的两种镜像变体，1911 年他亲手把这种“镜像分子”拆分开来——这一手几乎堵死了所有退路。1913 年，他因此获得诺贝尔化学奖。

这个故事值得记住两点。第一，配位理论不是靠一个天才实验一锤定音的，而是模型给出定量预言（几份氯能沉淀）、实验逐个排除对手，一路缠斗了二十年。第二，维尔纳那时根本不知道 d 电子是什么——电子刚被发现，量子力学还没出生。他靠颜色和沉淀这些宏观现象，倒推出了原子的空间排布。先摸清规律，再补机制，这是科学的常态，不丢人。
\end{zhengju}

\section{生命抢先把过渡金属用到了极致}

如果你觉得配位和催化只是工厂和宝石店的事，那就小看自然了。生命在三十多亿年前就把过渡金属的这三样特产全数征用。

先说你的血。血液呈红色，靠血红蛋白，而每个血红蛋白分子的正中心坐着一个配位化合物：铁离子被一个叫“血红素”的环状配体从四个方向抱住，剩下的位置留给氧气。\chem{Fe^{2+}} 与 \chem{O_2} 的结合讲究的是\textbf{恰到好处}：配体把结合力调得不太强也不太弱——在肺里氧气多就抓住，到了组织氧气少就松手。一氧化碳的可怕正在于此：它钻进同一个位置，抱得比氧气紧两百多倍，赖着不走，血红蛋白的车就被占满了。为什么非铁不可？因为铁的 d 电子和多变价态，正好胜任“可逆地抓氧又放氧”这份工作——这是主族金属给不了的特长。

再说钴。维生素 \chem{B_{12}} 是已知结构最复杂的维生素，中心坐着一个钴离子，同样被环状配体环抱。人体自己造不出它，靠肠道细菌和食物供给；缺了它，造血和神经都会出故障。还有植物：豆科植物根瘤里的固氮酶，用钼和铁搭成的催化中心，在常温常压下完成化肥厂需要高温高压才能做的固氮反应——顺便说一句，这个中心的精确工作原理，至今仍有细节没被完全搞清楚，科学在这里还没交卷。

给你算一笔账，感受下数量级：一个成年人全身大约含 \ud{4}{g} 铁，大部分在血红蛋白里。铁的摩尔质量约 \ud{56}{g\,mol^{-1}}，所以 $4/56 \approx 0.07$ mol，乘以每摩尔 \ud{6.02\times 10^{23}}{} 个原子：
\[0.07 \times 6.02\times 10^{23} \approx 4\times 10^{22}\ \text{个铁原子}\]
四万亿亿个铁原子，此刻正在你的血管里轮流抓氧、放氧。你每一次呼吸，都依赖过渡金属的配位化学正常运转。

\section{本章模型的边界在哪里}

老规矩，说完好用的部分，必须说它会怎么失效。

第一，“颜色来自 d 电子在房间间跃迁”这个解释对很多离子成立，但不是全部。高锰酸根 \chem{MnO_4^{-}} 的锰是 $+7$ 价，d 房间其实是空的，根本没有 d 电子可以跃迁——它的深紫色来自另一类过程：电子从氧配体整块“搬”向锰（电荷转移），吸收光的本领更强，颜色也更深。所以“过渡金属有色全靠 d 电子”是个好用的一阶近似，不是完整理论。

第二，“电子住房间、自旋像小磁针”从头到尾是模型。它解释趋势、预言例外时表现出色，但电子不是真的小磁针，磁畴对齐的深层原因要动用量子力学里更微妙的交换作用。遇到定量问题——比如精确预测某种合金的磁性——本章的图景只能指路，不能结算。

第三，催化的表面故事被大幅简化了。真实催化剂表面凹凸不平，起作用的往往只是少数“活性位点”；杂质、温度、颗粒大小的影响都可能颠覆结论。“铁能催化合成氨”是对的，但“为什么这一块铁比那块铁好用”，至今还是催化学的前沿问题。

\section{回到开场那四样东西}

现在对答案。生锈铁钉与不锈钢勺：同为铁家族，差在铬。铬抢先锈出一层致密透明的氧化膜，把水和氧挡在门外；普通铁的锈疏松吸水，只会引狼入室。

红铜丝：铜的 3d 房间几乎住满，但还没满，d 电子跃迁恰好吃掉蓝紫色的光，把红橙色反射回来，于是铜呈红棕色；银的相应能量差落在紫外区，可见光毫发无损，所以银是银白的。金属的颜色，同样是 d 电子的账本。

磁铁：铁、钴、镍的 3d 房间半满，未成对电子多，且相邻原子间容易让小磁针集体对齐——两个条件缺一不可，所以磁铁三兄弟在周期表上做邻居不是巧合。

蓝色胆矾：铜离子被水分子团团围住，水配体把 d 房间的能量差拉到恰好吃掉橙黄光的位置，剩下的蓝紫光进入你的眼睛。把水烤干，颜色退场；换上氨配体，颜色加深。那杯蓝色，是铜离子和它身边六个水房客共同的签名。

\section{这块街区还没逛完}

这一章我们回答了“过渡金属为什么颜色多、有磁性、会催化”，但有两个更大的问题被悄悄带了出来。其一：铁、钴、镍、铜这些元素当初是从哪儿来的？地壳里的铁只占地核之外的一小部分，你体内的那 4 克铁、电机里的钴、手机里的稀土，全都经历过一段超乎想象的旅程——那是下一章（第 16 章）要讲的、从恒星熔炉到你血管的故事。其二：铁催化合成氨要高温高压，你细胞里的酶却在体温下做得更好更快。蛋白质凭什么把过渡金属的本事再放大几个数量级？这个问题要等到第二册，等我们先弄明白碳怎样搭出生命的大分子，再来揭晓。

\begin{tiaozhan}
为什么 d 电子跃迁的能量差“恰好”落在可见光范围，而不是深紫外或红外？可以粗略理解为两股力量的平衡：d 电子离核比 s、p 电子远不了太多，房间间的分裂能由配体电子对的排斥决定，量级在 \ud{100}{} 到 \ud{400}{kJ\,mol^{-1}} 之间。把一份光的能量 $E = hc/\lambda$ 换算成每摩尔：波长 $\lambda$ 越短能量越高。代入可见光的边界 $\lambda \approx 400$ nm 和 $700$ nm，对应的正是约 \ud{300}{} 到 \ud{170}{kJ\,mol^{-1}}——和配体场的分裂能几乎重叠。这不是巧合设计，而是电磁相互作用强度在原子尺度上的自然结果：若分裂能大十倍，所有配合物都会无色，若小十倍，它们全都会黑得像墨。进一步的问题——为什么五种 d 房间在八面体配体包围下会分裂成“三间低、两间高”——需要画出 d 房间的空间朝向和配体的方位，那是大学配位场理论的开篇。跳过这一段，丝毫不影响你理解本章主线。
\end{tiaozhan}

\section*{练一练}

\begin{sikaoti}
\item 画出 \chem{Fe^{2+}} 在水中配位的示意图：中心一个铁离子，周围画六个水分子（用水滴形状表示，标出氧端朝向铁）。图旁注明：哪些部分属于“内界”，为什么这个整体会有颜色，而一杯食盐水没有。
\item 铁在潮湿的空气中生锈，在干燥的沙漠里却保存千年。用“反应需要哪些粒子同时到场”的语言，画一幅物质流图：水、氧气、铁各自扮演什么角色？据此解释为什么船底水线附近锈得最严重。
\item 画一条能量曲线图：横轴是反应进程，纵轴是能量，画出反应物、产物和中间的能量山。再用虚线画出“加入催化剂后”的另一条路径。在图上标出：哪一段高度代表活化能？催化剂改变了什么、没有改变什么？
\item 红宝石和蓝宝石的骨架是同一种无色矿物氧化铝，差别只在万分之几的杂质。用本章的配位与 d 电子图像，向一位没读过本章的朋友解释：为什么混入不同的过渡金属杂质，石头会呈现不同颜色？杂质浓度那么低，为什么颜色却那么明显？
\item 有人宣称：“不锈钢是绝对不会生锈的金属，因为它不含会锈的成分。”找出这句话里的两处错误，并用铬和氧化膜的语言写出正确版本。
\item 查一查你家厨房的用品（锅、勺、水龙头、易拉罐），列出哪些用了过渡金属，并猜一猜各自利用了本章讲的哪一条性质（颜色、磁性、催化、致密氧化膜、导热等）。哪一条你猜不出来？记下来，看看后面的章节能不能回答。
\end{sikaoti}
